\documentclass[11pt]{article}

\usepackage{ods_etsinf}

\setlength{\parindent}{1.5em}
\setlength{\parskip}{0.8em}

\begin{document}

\phantom{x}

\vspace{1ex}

\section*{ANEXO}

\vspace{2ex}

OBJETIVOS DE DESARROLLO SOSTENIBLE

\vspace{4ex}

Grado de relaci\'on del trabajo con los Objetivos de Desarrollo Sostenible (ODS).

\vspace{2ex}

\begin{tabular}{|l|c|c|c|c|}\hline
\textbf{Objetivos de Desarrollo Sostenible} & \textbf{Alto} & \textbf{Medio} & \textbf{Bajo} & \textbf{No} \\
& & & & \textbf{procede} \\ \hline
ODS 1.  \textbf{Fin de la pobreza.}                            & & & & X \\ \hline
ODS 2.  \textbf{Hambre cero.}                                  & & & & X \\ \hline
ODS 3.  \textbf{Salud y bienestar.}                            & & & & X \\ \hline
ODS 4.  \textbf{Educaci\'on de calidad.}                       & X & & & \\ \hline
ODS 5.  \textbf{Igualdad de g\'enero.}                         & & & & X \\ \hline
ODS 6.  \textbf{Agua limpia y saneamiento.}                    & & & & X \\ \hline
ODS 7.  \textbf{Energ\'{\i}a asequible y no contaminante.}     & & & X & \\ \hline
ODS 8.  \textbf{Trabajo decente y crecimiento econ\'omico.}    & & & & X \\ \hline
ODS 9.  \textbf{Industria, innovaci\'on e infraestructuras.}   & X & & & \\ \hline
ODS 10. \textbf{Reducci\'on de las desigualdades.}             & & & X & \\ \hline
ODS 11. \textbf{Ciudades y comunidades sostenibles.}           & & & & X \\ \hline
ODS 12. \textbf{Producci\'on y consumo responsables.}          & & X & & \\ \hline
ODS 13. \textbf{Acci\'on por el clima.}                        & & & X & \\ \hline
ODS 14. \textbf{Vida submarina.}                               & & & & X \\ \hline
ODS 15. \textbf{Vida de ecosistemas terrestres.}               & & & & X \\ \hline
ODS 16. \textbf{Paz, justicia e instituciones s\'olidas.}      & & & & X \\ \hline
ODS 17. \textbf{Alianzas para lograr objetivos.}               & & & X & \\ \hline
\end{tabular}

\newpage

\vspace*{0.5cm}

\noindent\textbf{Reflexi\'on sobre la relaci\'on del TFG/TFM con los ODS y con el/los ODS m\'as relacionados.}


Este TFG, centrado en la transcripci\'on autom\'atica de piano polif\'onico mediante modelos de aprendizaje profundo, presenta una relaci\'on especialmente significativa con el ODS 9, relativo a la industria, la innovaci\'on y las infraestructuras. El trabajo no se limita al uso de herramientas ya existentes, sino que aborda la reproducci\'on, configuraci\'on, entrenamiento, evaluaci\'on y comparaci\'on de arquitecturas avanzadas como MFT-FiLM y Transkun sobre el conjunto MAESTRO v3.0.0. Esta aproximaci\'on contribuye al desarrollo de soluciones innovadoras en el \'ambito del procesamiento de se\~nales musicales y de la inteligencia artificial aplicada, reforzando adem\'as la importancia de la reproducibilidad experimental como parte esencial de una infraestructura cient\'ifica y tecnol\'ogica fiable.

La relaci\'on con el ODS 4, educaci\'on de calidad, tambi\'en resulta relevante. La transcripci\'on autom\'atica de m\'usica puede facilitar el acceso a partituras, an\'alisis musicales y materiales docentes generados a partir de grabaciones reales. Este tipo de tecnolog\'ia tiene un potencial evidente en contextos de formaci\'on musical, tanto para estudiantes como para docentes, ya que permite estudiar interpretaciones, comparar ejecuciones y analizar pasajes complejos de manera m\'as sistem\'atica. En particular, los modelos evaluados permiten convertir audio de piano en representaciones simb\'olicas, lo que puede apoyar el aprendizaje musical y favorecer nuevas formas de interacci\'on entre interpretaci\'on, an\'alisis y docencia.

Asimismo, el trabajo guarda relaci\'on con el ODS 12, producci\'on y consumo responsables. Aunque el entrenamiento de modelos neuronales requiere recursos computacionales significativos, el proyecto incorpora una reflexi\'on pr\'actica sobre el uso responsable de dichos recursos: configuraci\'on cuidadosa del entorno, ejecuci\'on mediante SLURM, reutilizaci\'on de \textit{checkpoints}, comparaci\'on entre modelos de distinto tama\~no y an\'alisis de eficiencia. El estudio de MFT-FiLM, un modelo competitivo con un n\'umero reducido de par\'ametros, resulta especialmente pertinente desde esta perspectiva, ya que muestra que no siempre es necesario recurrir a arquitecturas de gran escala para obtener resultados s\'olidos.

De forma m\'as indirecta, el trabajo se relaciona con el ODS 13, acci\'on por el clima, y con el ODS 7, energ\'ia asequible y no contaminante, en la medida en que la eficiencia computacional y la planificaci\'on de los entrenamientos permiten reducir ejecuciones innecesarias y optimizar el consumo energ\'etico asociado al c\'omputo de alto rendimiento. Tambi\'en existe una vinculaci\'on menor con el ODS 10, reducci\'on de las desigualdades, ya que las herramientas de transcripci\'on autom\'atica pueden disminuir barreras de acceso a la notaci\'on musical y favorecer que m\'as personas puedan estudiar, analizar o reutilizar contenido musical sin depender exclusivamente de transcripciones manuales costosas.

Finalmente, el ODS 17 aparece de forma transversal por el uso de repositorios, art\'iculos cient\'ificos, conjuntos de datos abiertos y metodolog\'ias comparables con trabajos previos. En conjunto, este TFG demuestra c\'omo un proyecto t\'ecnico de ingenier\'ia inform\'atica puede contribuir, desde la investigaci\'on aplicada y la reproducibilidad, al desarrollo de tecnolog\'ias m\'as accesibles, eficientes e innovadoras.
\end{document}